% Generated draft: requires mathematical and source review.
\documentclass[11pt]{article}
\usepackage[margin=25mm]{geometry}
\usepackage{fontspec}
\setmainfont{Latin Modern Roman}
\usepackage{amsmath,amssymb,mathtools,unicode-math}
\setmathfont{Latin Modern Math}
\usepackage{enumitem,needspace,xurl,hyperref}
\setlength{\parindent}{0pt}
\setlength{\parskip}{5pt}
\setlength{\emergencystretch}{3em}
\newcommand{\N}{\mathbb{N}}
\newcommand{\Z}{\mathbb{Z}}
\newcommand{\Q}{\mathbb{Q}}
\newcommand{\R}{\mathbb{R}}
\newcommand{\C}{\mathbb{C}}
\newcommand{\F}{\mathbb{F}}
\newcommand{\A}{\mathbb{A}}
\newcommand{\PP}{\mathbb{P}}
\newcommand{\E}{\mathbb{E}}
\newcommand{\eps}{\varepsilon}
\newcommand{\dd}{\,\mathrm{d}}
\begin{document}
\section*{UnsolvedMath ID: 30001437}
\textbf{Source problem code:} catalogue entry
\subsection{Definitions and mathematical statement}
The source concerns rigidity of Bernoulli actions. The formulation below uses the usual setting of countable discrete groups acting by probability-measure-preserving transformations on standard probability spaces. These conventions are interpretive qualifications: the record does not explicitly specify the measure-space category, preservation of probability, or countability of the comparison group.

A group \(\Gamma\) is ICC if every nonidentity element has an infinite conjugacy class. Let \(H\trianglelefteq\Gamma\) be an infinite normal subgroup. The inclusion \(H\leq\Gamma\) has relative property \((T)\) if every unitary representation \(\pi:\Gamma\to\mathcal U(\mathcal H)\) with almost invariant unit vectors has a nonzero \(H\)-invariant vector. Here almost invariance means that for every finite \(K\subseteq\Gamma\) and every \(\varepsilon>0\), there is a unit vector \(\xi\in\mathcal H\) such that \(\|\pi(g)\xi-\xi\|<\varepsilon\) for all \(g\in K\).

For a probability space \((X_0,\mu_0)\), the Bernoulli action is the shift on
\[
(X,\mu)=(X_0^{\Gamma},\mu_0^{\Gamma}),\qquad
(g\cdot x)_h=x_{g^{-1}h}\quad(g,h\in\Gamma).
\]
An action \(\Lambda\curvearrowright(Y,\nu)\) is essentially free if almost every point has trivial stabilizer, and ergodic if every invariant measurable set has measure zero or one. Its group-measure-space von Neumann algebra is the crossed product \(L^{\infty}(Y,\nu)\rtimes\Lambda\), generated by \(L^{\infty}(Y,\nu)\) and implementing unitaries \(u_\lambda\) satisfying \(u_\lambda f u_\lambda^*=f\circ\lambda^{-1}\).

\textbf{Question, under the stated interpretive conventions.} Let \(\Gamma\) be a countable ICC group admitting an infinite normal subgroup \(H\) for which \((\Gamma,H)\) has relative property \((T)\), and let \(\Gamma\curvearrowright(X,\mu)\) be a Bernoulli action. For every essentially free ergodic action \(\Lambda\curvearrowright(Y,\nu)\), does an isomorphism of von Neumann algebras
\[
L^{\infty}(X,\mu)\rtimes\Gamma
\cong
L^{\infty}(Y,\nu)\rtimes\Lambda
\]
imply the existence of a group isomorphism \(\delta:\Gamma\to\Lambda\) and a measure-preserving isomorphism \(\theta:(X,\mu)\to(Y,\nu)\), modulo null sets, such that
\[
\theta(g\cdot x)=\delta(g)\cdot\theta(x)
\qquad\text{for every }g\in\Gamma\text{ and almost every }x\in X?
\]
This specifies conjugacy through an isomorphism of the acting groups. The source does not expressly define conjugacy in this setting. It also does not specify nontriviality of the Bernoulli base or explicitly require the Bernoulli action itself to be essentially free; no such additional hypotheses have been silently inserted.

The supplied curation labels the record partially solved and attributes an affirmative result for the stated group class to Ioana's W-star-superrigidity theorem. It also describes an unrestricted extracted formulation as broader. That broader formulation is not separately supplied, so its scope cannot be reconstructed here. These are supplied literature qualifications, not independently verified conclusions.
\subsection{Short English statement}
For a Bernoulli action of a countable ICC group containing an infinite normal subgroup with relative property \((T)\), does its group-measure-space von Neumann algebra determine the action up to conjugacy through a group isomorphism? The probability-space conventions and some action hypotheses require clarification. The supplied dated curation attributes an affirmative result in this group class to Ioana while retaining a partially solved label for a reportedly broader question.
\subsection{Sources}
\begin{itemize}
\item Local UnsolvedMath registry: permanent ID \(30001437\), source problem code \texttt{catalogue entry}, title ``Conjugacy Rigidity from Group-Measure-Space Factors.'' Registry locations specified by the project are \texttt{data/unsolvedmath/problems.json} and the supplementary \texttt{catalog/problems.json}; this draft uses the supplied complete record without independently inspecting those files.
\item Supplied report citation: \emph{C*-Algebren} (2010), Oberwolfach Reports, DOI 10.4171/owr/2010/13, pp.\ 667--732. \url{https://doi.org/10.4171/owr/2010/13}. The record attributes its provenance to Oberwolfach workshop reports and open-problem discussions; no individual proposer is supplied.
\item Supplied literature source: ``Oberwolfach Report 13/2010.'' \url{https://doi.org/10.4171/owr/2010/13}.
\item Supplied literature source: Ioana, ``W-star-superrigidity for Bernoulli actions.'' \url{https://arxiv.org/abs/1002.4595}.
\item Literature assessment dated exactly \texttt{2026-08-21}: the supplied curation labels the problem partially solved and attributes an affirmative answer for Bernoulli actions of ICC groups with an infinite normal subgroup having relative property \((T)\) to Ioana, while describing unrestricted extracted wording as broader. This is supplied dated curation, not an independently verified literature assessment. The background expressly requests independent verification of time-sensitive status and cited claims.
\item The supplied statement status is \texttt{corrected\_verified}; its accompanying verification note reports only insertion of mathematical delimiters without changing mathematical content. This does not establish mathematical correctness or verify the cited theorem's hypotheses.
\end{itemize}
\end{document}
